\documentclass[11pt,letterpaper]{article}
\usepackage[T1]{fontenc}
\usepackage[utf8]{inputenc}
\usepackage{lmodern}
\usepackage[margin=0.9in,headheight=15pt]{geometry}
\usepackage{microtype}
\usepackage{amsmath,amssymb,amsthm}
\usepackage{booktabs,tabularx,array,longtable}
\usepackage[table]{xcolor}
\usepackage{listings}
\usepackage{enumitem}
\usepackage{titlesec}
\usepackage{needspace,etoolbox}
\usepackage{fancyhdr}
\usepackage{xurl}
\usepackage{hyperref}
\definecolor{ink}{HTML}{18334D}
\definecolor{muted}{HTML}{566879}
\definecolor{wash}{HTML}{F2F5F7}
\definecolor{linkblue}{HTML}{205F87}
\hypersetup{colorlinks=true,linkcolor=ink,urlcolor=linkblue,citecolor=linkblue,
  pdftitle={AsymptoticAnalysis: Incremental Audit of Logarithmic Certificate Precision},
  pdfauthor={OpenAI},pdfsubject={Pinned-source review, exact interval proofs, native witnesses and candidate repairs}}
\titleformat{\section}{\Large\sffamily\bfseries\color{ink}}{\thesection}{0.7em}{}
\titleformat{\subsection}{\large\sffamily\bfseries\color{ink}}{\thesubsection}{0.7em}{}
\titleformat{\subsubsection}{\normalsize\sffamily\bfseries\color{ink}}{\thesubsubsection}{0.7em}{}
\titlespacing*{\section}{0pt}{2.3ex plus .5ex}{1.2ex}
\titlespacing*{\subsection}{0pt}{1.8ex plus .5ex}{.8ex}
\setlength{\parindent}{0pt}
\setlength{\parskip}{.55em}
\setlength{\emergencystretch}{2em}
\setcounter{tocdepth}{2}
\pretocmd{\section}{\Needspace{6\baselineskip}}{}{}
\clubpenalty=10000
\widowpenalty=10000
\setlist{itemsep=2pt,topsep=4pt,leftmargin=1.5em}
\pagestyle{fancy}
\fancyhf{}
\fancyhead[L]{\small\sffamily\color{muted} AsymptoticAnalysis \; / \; Incremental source audit}
\fancyhead[R]{\small\sffamily\color{muted} 10 September 2026}
\fancyfoot[L]{\small\sffamily\color{muted} Snapshot efa1aeec4845}
\fancyfoot[R]{\small\sffamily\thepage}
\renewcommand{\headrulewidth}{0.3pt}
\lstset{basicstyle=\ttfamily\footnotesize,backgroundcolor=\color{wash},
  frame=single,rulecolor=\color{wash},framesep=6pt,xleftmargin=6pt,xrightmargin=6pt,
  columns=fullflexible,keepspaces=true,showstringspaces=false,breaklines=true,
  breakatwhitespace=false,tabsize=2,aboveskip=10pt,belowskip=10pt,
  upquote=true}
\newcommand{\code}[1]{\texttt{\detokenize{#1}}}
\newcommand{\src}[1]{\texttt{\small\detokenize{#1}}}
\newcommand{\R}{\mathbb R}
\newcommand{\Q}{\mathbb Q}
\newcommand{\I}{\mathcal I}
\newcommand{\width}{\operatorname{wid}}
\newcommand{\roundout}{\operatorname{out}}
\newtheorem{proposition}{Proposition}[section]
\newtheorem{lemma}[proposition]{Lemma}
\theoremstyle{remark}\newtheorem{remark}[proposition]{Remark}
\newcolumntype{Y}{>{\raggedright\arraybackslash}X}

\begin{document}
\begin{titlepage}
\vspace*{0.55in}
{\sffamily\color{muted}\large PINNED-SOURCE TECHNICAL REVIEW\par}
\vspace{0.35in}
{\sffamily\bfseries\color{ink}\Huge AsymptoticAnalysis\par}
\vspace{0.18in}
{\sffamily\color{ink}\LARGE Logarithmic certificate precision\par}
\vspace{0.22in}
{\large Two independent losses of relative accuracy,\par
native reproductions, proofs, and candidate repairs\par}
\vspace{0.55in}
{\renewcommand{\arraystretch}{1.22}
\begin{tabularx}{\textwidth}{@{}p{1.18in}Y@{}}
\toprule
Repository & \url{https://github.com/VladimirReshetnikov/Asymptotic} \\
Revision & \path{efa1aeec4845a9c35e140963a0333d0c9ec33b05} \\
Review date & 10 September 2026 \\
Native evidence & Wolfram Language 15.0.0, Linux x86-64; downloaded standalone \\
Other evidence & Exact-rational Python models; 24 passing test methods \\
Mathics boundary & Source analysis and supplied regressions; no Mathics execution \\
\bottomrule
\end{tabularx}}\par
\vspace{0.16in}
\textbf{Principal result.} Two distinct implementation choices discard relative
accuracy around a logarithm's unit argument. The baseline safely refuses three
exact-root certificate requests. Two small, independently tested changes make
all three requests certify at the same enclosure order. No false certificate
or incorrect expansion is alleged by these findings.

\textbf{Novelty boundary.} This is an incremental review after six existing
review waves. Known findings and their general roadmaps are excluded from the
new work. The retained contribution is deliberately narrower than a claim of
exhaustive verification of every repository file.
\vfill
{\small\color{muted}The archive contains the article source, candidate emitter,
reproduction specifications, independent models, selected native observations,
and explicit evidence limitations. Candidate code is not an upstream change.}
\end{titlepage}

\begin{abstract}
This review investigates the current AsymptoticAnalysis implementation against
its Wolfram and Mathics compatibility objective while screening out obligations
already recorded in the repository's external reviews. Two new shortcomings
remain after that screening. First, a positive rational just below one is
reduced to a logarithm near two minus a separately enclosed logarithm of two;
the resulting interval has an absolute uncertainty essentially unrelated to
the tiny logarithm being computed. Second, an exactly computed rational affine
argument of a logarithm is rounded near one before the logarithm is evaluated,
losing the small displacement that determines the answer. These mechanisms
are distinct: a four-variant, three-witness native experiment separates them.

The proposed repairs use reciprocal reduction below one and preserve exact
affine endpoints until after taking the logarithm. Their containment arguments
are proved, their selected public consequences are reproduced on Wolfram
15.0.0, and independent rational models pass 24 test methods. The article also
develops a narrowly scoped, positivity-guarded logarithmic-ratio residual
prototype. Mathics execution, canonical-source integration, a complete MUnit
run, and performance benchmarking remain outside the executed evidence.
\end{abstract}
\tableofcontents
\clearpage

\section{Executive assessment}

\subsection{What should change}
The highest-value new work found in this audit is concentrated in
\src{src/Kernel/InverseCertificates.wl}. The certificate engine is already
careful about outward rounding, real domains, derivative separation, and
existence evidence. The problems below are not failures of those safety
checks. They make otherwise straightforward requests consume unnecessary
precision or fail to produce a certificate at the requested resource
settings. The inspected source is the pinned revision, not an inferred
current branch state.\cite{source,certsource}

\begin{table}[htbp]
\centering\small
\begin{tabularx}{\textwidth}{@{}p{.5in}Yp{1.46in}@{}}
\toprule
ID & New obligation & Evidence and priority \\
\midrule
F01 & Preserve relative precision when evaluating an exact rational logarithm
just below one. Replace cancellation of two near-equal log enclosures by a
single reciprocal reduction. & Native reproduction and repair; functional
robustness/performance priority. \\
F02 & Do not round an exact affine logarithm argument near one before applying
the logarithm. Pass its exact range to the logarithm primitive. & Native
reproduction and repair; functional robustness/performance priority. \\
E01 & Explore a positive logarithmic-ratio residual that retains an exact
zero at a known rational center. & New scoped development proposal;
primitive checked, dispatcher integration unrun. \\
\bottomrule
\end{tabularx}
\caption{Retained deltas. F01 and F02 are safe-refusal shortcomings, not
unsound certificates. E01 is not counted as a demonstrated defect.}
\end{table}

For $\delta=2^{-200}$, the original low-order point enclosure of
$\log(1-\delta)$ straddles zero with endpoint magnitude approximately
$1.85\times10^{-3}$, although the logarithm has magnitude about
$6.22\times10^{-61}$. Separately, evaluating $\log(1+x)$ at the exact point
$x=\delta$ produces an enclosure approximately $[0,7.11\times10^{-15}]$.
Both are valid enclosures; both are too wide for the deliberately small
verification intervals in the public witnesses.

The combined candidate certifies all three native witnesses at
\code{"EnclosureOrder" -> 2} and \code{"MaxRefinements" -> 0}. In every
successful row of the controlled comparison, the returned root interval
contains the known exact root. The corresponding bounds on error of the
\emph{supplied center} are about $4.42\times10^{-75}$ or
$2.21\times10^{-75}$. They are not automatically error bounds for the stored
finite asymptotic expression.

\subsection{What this result does not establish}
An enclosure-order option is a resource choice, not a promise that every
problem must certify at that order. Thus a baseline refusal alone would be
weak evidence of a bug. The stronger evidence is the identified loss mechanism,
a family showing its scaling, a containment-preserving repair, and the
same-resource native comparison. These justify treating the behavior as an
avoidable numerical and coverage shortcoming without portraying a safe
failure as an incorrect theorem.

The successful candidate observations do not establish complete Wolfram
compatibility, complete Mathics compatibility, or general superiority in every
logarithm calculation. No default-setting benchmark was run. In particular,
the report does not claim that ordinary adaptive defaults always fail for
these moderate-size examples: increasing arithmetic work can recover some
lost information. The issue is that substantially more work is then demanded
for an avoidable representation choice.

\section{Scope, provenance, and nonduplication}
\subsection{Pinned source and execution boundary}
The reviewed commit is
\path{efa1aeec4845a9c35e140963a0333d0c9ec33b05}, committed on
10 September 2026 at 19:09:04 UTC. Its root standalone package contains
692,381 bytes. The remote native experiments downloaded this exact standalone
to a temporary file and loaded it into a Wolfram Language kernel identifying
itself as \code{15.0.0 for Linux x86 (64-bit) (May 6, 2026)}.
The four-variant comparison reloaded temporary standalone copies within one
kernel; it is not four independent environments. Other focused native probes
were separate calls.\cite{source}

A local repository clone and a local Mathics installation could not be
obtained because the local environment's network access failed. The GitHub
connector supplied source text, and the remote Wolfram evaluator supplied
selected package executions. The remote service was intermittent: failed
service calls are not counted as package failures, successful tests, or
completed corpus scans.

The following distinctions apply throughout:
\begin{description}[style=nextline]
\item[Native observation]
A reported value came from a successful Wolfram call loading the pinned
standalone or an explicitly identified temporary candidate.
\item[Independent model]
Python's \code{Fraction} arithmetic implements the small enclosure mechanism.
It does not execute the repository and is not a Mathics simulation claim.
\item[Proof]
A mathematical argument establishes an enclosure or scaling property under
stated exact-real hypotheses. It is not a machine-checked formalization.
\item[Regression specification]
A supplied WL/WLT file records desired behavior. Its existence is not evidence
that the complete file or MUnit suite was executed.
\end{description}

No upstream files were modified. No canonical modular build, standalone
freshness check, full package suite, complete new portable suite, or
wall-clock performance benchmark was run. The code in the archive is a
candidate and a reproduction aid, not a replacement release.

\subsection{Breadth of inspection}
The source inspection extended beyond the two retained findings. It covered
public request and result definitions; substantial parts of the forward
power-log jet arithmetic; the fixed-parameter special-function adapter; exact
special-function identity normalization; real-domain guards; envelope
arithmetic; selected Gamma, Dirichlet, and inverse-check paths; Mathics
assumption, numerical, list, and compatibility adapters; and most of the
certificate module. The fixed-parameter and identity modules were inspected
particularly for lost branches, incorrect offsets, and finite identities
mistaken for asymptotic truncations. No additional new defect in those paths
was established by the performed analysis.\cite{kernel,specialparam,specialident}

This is not a claim to have read every module line by line or to have
verified all formulas reachable through native special functions. A broad
search followed by a small retained set is preferable to padding the report
with already-known or unverified defects. The useful deliverable here is a
precise incremental contribution with runnable evidence.

\subsection{How existing findings were excluded}
The repository index records 45 retained packages across six review waves,
with ten retired packages. Its count of 231 entries for waves 1--4 is a count
of attributed report entries before consolidation, not 231 distinct current
bugs. This audit used all six wave indexes and substantial portions of the
maintained implementation register to screen candidate findings.\cite{reviews,status,wavefive,wavesix}

A targeted additional text screen fetched twelve retained wave-2 TeX files
and searched for the certificate logarithm helper names and near-one/log1p
terms. That completed scan had no fetch failures and no matches. Larger
attempted scans failed and are not claimed as completed. Not every retained
article or PDF was individually reread. Consequently the defensible novelty
claim is that these two mechanisms are absent from the inspected maintained
crosswalks and targeted material, not a logically exhaustive proof that no
sentence in the entire archive could be related.

\begin{table}[htbp]
\centering\small
\begin{tabularx}{\textwidth}{@{}p{1.0in}YY@{}}
\toprule
Existing item & Previously recorded concern & Why the retained delta differs \\
\midrule
C19 & Certificate accuracy planning, progress, tolerances, and refinement. &
F01 and F02 occur inside a single exact arithmetic attempt. They remain
separable with adaptation disabled. \\
C20 & Preserve cancellation when finding an affine range, especially with a
large translated source coordinate. & F02 starts after that exact range has
already been computed correctly. The later rounding before a logarithm
loses a different small quantity. \\
P08 & General profiling of proof and arithmetic work. & The new report gives
a specific operation, a parameterized loss mechanism, and executed local
repairs, not another generic profiling recommendation. \\
Wave 6 & Mathics logarithm splitting authorized by rounded sign tests;
interval-power geometry and budget behavior. & The new witnesses use the
shared rational logarithm enclosure with positive exact inputs. No
machine-sign split or interval power is involved. \\
\bottomrule
\end{tabularx}
\caption{Nearest existing obligations are acknowledged, not restated as new
findings. The archive contains a separate novelty crosswalk.}
\end{table}

The existing generic capability, native-backend, multiscale, certification,
release, and provenance roadmaps are not repeated as new recommendations.
The implementation plan below is restricted to the two repairs and their
centered-residual extension.

\section{The certificate arithmetic relevant to the findings}
\subsection{Exact intervals and directed significant-bit rounding}
Let an interval $[a,b]$ with rational endpoints represent a proved real
inclusion. For a nonzero rational $q$, the code chooses
\[
 e=\lfloor\log_2|q|\rfloor,\qquad h=2^{e-B+1},
\]
where $B$ is the significant-bit budget. It rounds downward with
$h\lfloor q/h\rfloor$ and upward with $h\lceil q/h\rceil$.
Addition and multiplication compute exact rational endpoint expressions,
then round outward. No floating-point sign test is used in this
primitive.\cite{certsource}

This arrangement protects inclusion. It does not retain dependence between
separately enclosed occurrences of almost the same quantity. Nor does exact
rational evaluation prior to one rounding step guarantee that a later
ill-conditioned transformation will preserve useful relative accuracy.
Those are the respective sources of F01 and F02.

The public driver sets
\begin{equation}\label{eq:bits}
 B=4(n+10),
\end{equation}
where $n$ is \code{"EnclosureOrder"}. It accepts $2\le n\le2000$.
The main witnesses intentionally use $n=2$, hence $B=48$, and zero
refinements. They give an exact rational center strictly inside an exact
rational interval, so the numerical seed path is not the cause of failure.

\subsection{The existing logarithm formula}
For $1\le q\le2$, set
\[
 u=\frac{q-1}{q+1},\qquad 0\le u\le\frac13.
\]
The classical identity
\begin{equation}\label{eq:atanh}
 \log q=2\sum_{k=0}^{\infty}\frac{u^{2k+1}}{2k+1}
\end{equation}
is the logarithmic form of the inverse hyperbolic tangent series.\cite{dlmf}
The code retains $n$ terms and bounds the remaining positive tail:
\begin{align}
 S_n(q)&=2\sum_{k=0}^{n-1}\frac{u^{2k+1}}{2k+1},\label{eq:sum}\\
 T_n(q)&=\frac{2u^{2n+1}}{(2n+1)(1-u^2)}.\label{eq:tail}
\end{align}
Indeed, for $k\ge n$ the denominator $2k+1$ is at least $2n+1$, so the
remaining series is bounded by a geometric series. Thus
\begin{equation}\label{eq:unitbound}
 S_n(q)\le\log q\le S_n(q)+T_n(q).
\end{equation}
The final outward rounding preserves this enclosure.

For a general positive rational, the original point routine takes
$e=\lfloor\log_2q\rfloor$, $z=q/2^e\in[1,2)$, and combines independent
enclosures using
\begin{equation}\label{eq:reduction}
 \log q=\log z+e\log2.
\end{equation}
The identity is correct. Its interval evaluation is the problem when the
answer is a tiny difference of two quantities close to $\log2$.

\subsection{Why a wide residual prevents certification}
Suppose the verification interval is $I=[a,b]$, the candidate center is
$c\in(a,b)$, the stored forward function is continuously differentiable on
$I$, and the derivative has a fixed sign with $|f'|\ge\mu>0$ there. If the
residual enclosure proves $|f(c)-y|\le\varepsilon$, the natural residual
radius is $\varepsilon/\mu$.

When the resulting symmetric bracket is contained in $I$, continuity and
the derivative sign establish a root in that bracket. To see this in the
increasing case, use
\[
 f(c-r)-y\le f(c)-y-\mu r\le0,\qquad
 f(c+r)-y\ge f(c)-y+\mu r\ge0,
\]
with $r=\varepsilon/\mu$. The decreasing case is analogous. The code can
also establish existence through separately enclosed endpoint residual
signs. After existence is known, its signed derivative enclosure sharpens
the root interval.\cite{certsource}

A residual bound of size $10^{-3}$ or $10^{-15}$ cannot establish the desired
local bracket inside an interval of width about $10^{-60}$ when the
endpoint residuals are similarly overestimated. Refusing the certificate is
then correct. Improving the arithmetic that supplied those bounds, rather
than weakening existence or derivative checks, is the appropriate repair.

\section{F01: near-one cancellation in the point logarithm}
\subsection{Source mechanism}
The relevant routine is \code{certLogPoint}, at canonical source lines
61--72. After rejecting nonpositive arguments and returning an exact zero
for one, it performs the binary range reduction in
\eqref{eq:reduction}. There is no special treatment of a rational below
one.\cite{certsource}

Let $q=1-\delta$ with $0<\delta<1/2$. Then $e=-1$, so the implemented
calculation is
\begin{equation}\label{eq:badreduce}
 \log(1-\delta)=\log(2-2\delta)-\log2.
\end{equation}
The intervals for the two logarithms are treated as independent, although
they approach the same value as $\delta\to0^+$.

\begin{proposition}[Fixed-order loss]\label{prop:loss}
Before outward rounding, the interval produced by the two unit-log
bounds in \eqref{eq:badreduce} tends to
\[
 [-T_n(2),T_n(2)]\quad\text{as }\delta\to0^+,
\]
where
\begin{equation}\label{eq:tailtwo}
 T_n(2)=\frac{3^{1-2n}}{4(2n+1)}>0.
\end{equation}
For every fixed $n$, sufficiently small positive $\delta$ therefore makes
this enclosure straddle zero. Outward rounding cannot repair the sign loss.
\end{proposition}
\begin{proof}
The lower and upper endpoints of the unrounded subtraction are
\[
 S_n(2-2\delta)-S_n(2)-T_n(2)
\]
and
\[
 S_n(2-2\delta)+T_n(2-2\delta)-S_n(2).
\]
Continuity of these finite rational expressions near two gives the stated
limits. Substituting $u=1/3$ into \eqref{eq:tail} gives
\eqref{eq:tailtwo}. Since the limiting lower endpoint is strictly negative
and the limiting upper endpoint strictly positive, the crossing persists
for sufficiently small $\delta$. Enlarging an interval by outward rounding
cannot remove that crossing.
\end{proof}

For $n=2$, $T_2(2)=1/540$. The problem is not an algebraic identity error or
inexact input: it remains in exact rational arithmetic with valid explicit
tail bounds. At $q=1$ itself the exact-zero branch is used; the precision
behavior is sharply different immediately below it.

\subsection{Native witness A}
The following is the mathematical core of the executed request. Load the
pinned standalone in a separate input expression first; the archive's
portable probe uses fully qualified package symbols.
\begin{lstlisting}
d = 2^-200;
s = AsymptoticInverse[Log[x], {x, 1}, {y, 3},
      Direction -> "FromBelow"];
r = InverseCertificate[s, Log[1-d],
      "Interval" -> {1-2d, 1-d/2}, "Center" -> 1-d,
      "EnclosureOrder" -> 2, "MaxRefinements" -> 0];
\end{lstlisting}
The exact root is $x_*=1-d$. The interval lies strictly below the expansion
point one and contains that root in its interior. On the interval,
$f'(x)=1/x>1$. The native failure reported a derivative lower bound of one,
so failure to separate the derivative from zero is not involved.

The baseline returned
\code{Failure["ResidualBracketOutsideInterval", ...]}. A separate call to
the point logarithm primitive at the same order displayed approximately
\[
 [-0.00185185185185,\;0.00185185185185].
\]
That primitive's full exact interval was not returned in the successful
native transcript. The independent rational model gives the exact endpoints
\[
 \pm\frac{65156244609}{35184372088832},
\]
whose numerical values agree with the displayed native result. These exact
endpoints are model evidence, not an invented exact native transcript.

The root was supplied exactly, and the symbolic target was \code{Log[1-d]}.
No reference-root numerical solver, approximate target, or unproved
asymptotic error constant is needed to explain the failure.

\subsection{Candidate repair and inclusion proof}
Insert the following line immediately after the existing exact-one case,
with the existing positivity check retained:
\begin{lstlisting}
If[q < 1,
  Return[certNeg[certLogPoint[1/q, ctx]], Module]];
\end{lstlisting}
All arguments to the added branch are positive exact rationals. If $q<1$,
then $1/q>1$, so recursion occurs at most once. Arguments at least one
continue through the original branch, including the exact result at one.

\begin{proposition}[Reciprocal repair]
If the original positive-input logarithm routine returns a containing
interval $[L,U]$ at $1/q$, then the candidate interval $[-U,-L]$ contains
$\log q$. For $q=1-\delta$ and $0<\delta<1/2$, the new unit-series variable is
\begin{equation}\label{eq:smallu}
 u'=\frac{\delta}{2-\delta}.
\end{equation}
At fixed $n$, its unrounded tail is $O(\delta^{2n+1})$.
\end{proposition}
\begin{proof}
The positive-real identity $\log q=-\log(1/q)$ reverses the interval
endpoints and proves inclusion. Since $1<1/(1-\delta)<2$, the binary
exponent of the reciprocal is zero. Substitution into $(z-1)/(z+1)$ gives
\eqref{eq:smallu}. Equation \eqref{eq:tail} then gives the stated tail order.
\end{proof}

Because $|\log(1-\delta)|\sim\delta$, the relative tail is
$O(\delta^{2n})$. Outward rounding of the resulting small logarithm has
absolute scale $O(\delta 2^{-B})$, rather than the scale of a number near
$\log2$. This explains why the repair helps even though it does not increase
the significant-bit budget.

The proof establishes containment, not a claim that the new interval is
narrower than the original for every positive rational far from one. Nor
does it assert constant bit complexity: representing an input such as
$1-2^{-m}$ already requires a number of bits growing with $m$.

\subsection{Executed effect}
The patch anchor matched the downloaded standalone exactly once. Loading a
temporary copy with the inserted line changed witness A to a successful
certificate. The returned root enclosure contained $1-2^{-200}$, and its
bound on error of the supplied center was approximately
$4.42\times10^{-75}$. The primitive's observed width divided by $2^{-200}$
was about $7.11\times10^{-15}$; the exact independent model value is $2^{-47}$.

This is a controlled improvement with the same input interval, center,
enclosure order, and zero refinement budget. It is not an acceptance run of
all certificates or a measurement of end-to-end speed.

\section{F02: rounding an affine argument before its logarithm}
\subsection{Where correct exact work is discarded}
The package already contains \code{certAffinePair} and
\code{certAffineRange}. They recognize rational affine trees built from
constants, the source symbol, \code{Plus}, and \code{Times}, and compute an
exact range. That existing mechanism is not missing and is not proposed
again here.\cite{certsource,status}

However, \code{certEnclose} rounds that exact affine range before returning
it. The fallback in \code{certLogExpression} then applies
\code{certLog} to the rounded argument. The loss occurs in the sequence
\begin{equation}\label{eq:premature}
 \text{exact affine range}
 \;\longrightarrow\;\text{rounded range near 1}
 \;\longrightarrow\;\text{logarithm enclosure}.
\end{equation}
The relevant source is the logarithm-expression helper at lines 80--94
and the affine enclosure path at lines 109--130.

For $x=\delta=2^{-200}$, the exact value of $1+x$ is available. At $B=48$
its significant-bit grid near one has spacing $2^{-47}$, so outward rounding
turns the point into $[1,1+2^{-47}]$. Taking a logarithm of that interval must
include zero. Its width is governed by the precision near one, not by the
size of $\log(1+\delta)$.

\begin{proposition}[Information lost before a logarithm]\label{prop:affineloss}
Let $B\ge2$, $h=2^{1-B}$, and $0<\delta<h$. The significant-bit outward
rounding of the singleton $1+\delta$ is $[1,1+h]$. Every sound enclosure of
the logarithm over that rounded interval contains both $0$ and
$\log(1+h)$. Increasing only the accuracy of the subsequent logarithm
calculation cannot recover the original singleton's information.
\end{proposition}
\begin{proof}
All values in $[1,2)$ use exponent zero and grid $h$. For $0<\delta<h$, the
downward grid neighbor is one and the upward grid neighbor is $1+h$.
The logarithm is strictly increasing on the positive real axis, and its
exact image of that rounded interval is $[0,\log(1+h)]$. Any sound enclosing
interval must contain this image, irrespective of how accurately its
endpoints are computed.
\end{proof}

This is distinct from F01. For an exact point just above one, the original
point logarithm has binary exponent zero and does not subtract a logarithm
of two. The loss in F02 happens before that otherwise favorable primitive
receives its input.

\subsection{Native witness B}
\begin{lstlisting}
d = 2^-200;
s = AsymptoticInverse[Log[1+x], {x, 0}, {y, 3}];
r = InverseCertificate[s, Log[1+d],
      "Interval" -> {d/2, 2d}, "Center" -> d,
      "EnclosureOrder" -> 2, "MaxRefinements" -> 0];
\end{lstlisting}
Here the exact inverse is $e^y-1$, the supplied root is $x_*=d$, and the
package's retained finite expression was \code{y + y^2/2}. On the given
positive interval, $f'(x)=1/(1+x)$ is positive and close to one.

The baseline again returned
\code{Failure["ResidualBracketOutsideInterval", ...]}. Its recorded
residual enclosure was
\[
 [-2^{-200},\;2^{-47}],
\]
and its derivative lower bound was the positive rational
\[
 \frac{140737488355327}{140737488355328}.
\]
A primitive call displayed the logarithm enclosure
$[0,7.105427357601\times10^{-15}]$ at the exact singleton source point.
The interval is conservative, not incorrect, but its uncertainty is much
larger than the verification interval.

A separate decimal display of \code{N[Log[1+2^-200],12]} in that native batch
emitted \code{N::meprec}. That display is not used as a twelve-digit oracle
in this article. The package received an exact symbolic target, and the
independent numerical controls use \code{mpmath.log1p} at sufficient
precision. Confusing a display warning with an exact-certificate failure
would obscure the actual mechanism.

\subsection{Candidate repair}
Before the existing special cases in \code{certLogExpression}, add:
\begin{lstlisting}
If[MemberQ[{Plus, Times}, Head[argument]],
  candidate = certAffineRange[argument, x, interval];
  If[ListQ[candidate],
    Return[certLog[candidate, ctx], Module]]];
\end{lstlisting}
The existing \code{candidate} local is reused. The change does not ask for a
new general symbolic solver. It consumes the exact range already supported
by the package and delays rounding until the logarithm primitive has
converted the near-unit argument into the small answer.

\begin{proposition}[Exact affine logarithm range]
Let $p(x)=ax+b$ with $a,b\in\Q$, let $I=[\ell,r]$ have rational endpoints,
and let the exact range $J=p(I)$ be strictly positive. If the point-log
routine encloses the logarithm of each endpoint of $J$, choosing its lower
bound at $\min J$ and upper bound at $\max J$ encloses $\log p(x)$ for every
$x\in I$.
\end{proposition}
\begin{proof}
An affine function attains its extrema at interval endpoints; sorting
$a\ell+b$ and $ar+b$ therefore gives its exact range, including when $a<0$.
The logarithm is increasing on $(0,\infty)$. Its range is consequently
$[\log\min J,\log\max J]$, which is enclosed by the specified point bounds.
The existing \code{certLog} check rejects a nonpositive lower endpoint,
so the repair does not discard the required domain condition.
\end{proof}

For a nonaffine argument, the recognizer returns failure and the original
fallback remains available. For a positive product whose individual factors
are negative, an exact affine range is sufficient: the repair does not need
to assert real logarithms of those individual negative factors. This is not
an unrestricted product-logarithm rewrite.

\subsection{Executed effect and costs}
The F02 anchor also matched exactly once in the full standalone. With this
repair alone, witness B certified, its exact root was contained, and the
center error bound was approximately $2.21\times10^{-75}$. The primitive's
width divided by $2^{-200}$ was about $3.55\times10^{-15}$, with exact
independent model value $2^{-48}$.

For admitted affine expressions, exact endpoint construction was already
performed by the old path; the candidate avoids throwing away that
information before the logarithm. For a nonaffine \code{Plus} or
\code{Times} argument, the new dispatch can perform an extra unsuccessful
affine probe before the old path. Thus the repair is not advertised as a
universal speedup. A future implementation should measure this specific
additional work while retaining the mathematical invariant, rather than
changing unrelated affine recognizer architecture as part of this patch.

\section{Independent causes: the controlled native comparison}
\subsection{A translated negative-side witness}
The third witness is
\begin{lstlisting}
d = 2^-200;
s = AsymptoticInverse[Log[1+x], {x, 0}, {y, 3},
      Direction -> "FromBelow"];
r = InverseCertificate[s, Log[1-d],
      "Interval" -> {-2d, -d/2}, "Center" -> -d,
      "EnclosureOrder" -> 2, "MaxRefinements" -> 0];
\end{lstlisting}
Call this witness C. It is the same equation as A after the exact coordinate
translation $x_A=1+x_C$. The source interval and center are translated with
it. However, the expression-tree path now computes the affine argument
$1+x_C$, so both early rounding and below-one reduction can matter.

This is an especially useful regression because it does not rely on comparing
unrelated mathematical functions. A source coordinate change should not
manufacture orders of magnitude of avoidable proof uncertainty.

\subsection{Four full standalone variants}
The controlled experiment loaded the baseline and three temporary source
variants: reciprocal only, exact-affine logarithm only, and both. It evaluated
A, B, and C under each variant. The complete matrix is stored in
\src{evidence/native-matrix.json}; it is a selected transcription of a
successful remote response, not a raw MUnit receipt.

\begin{table}[htbp]
\centering
\begin{tabular}{@{}lccc@{}}
\toprule
Source variant & A: $\log x$ below 1 & B: $\log(1+x)$, $x>0$ & C: $\log(1+x)$, $x<0$ \\
\midrule
Baseline & Refused & Refused & Refused \\
F01 only & Certified & Refused & Refused \\
F02 only & Refused & Certified & Refused \\
F01 and F02 & Certified & Certified & Certified \\
\bottomrule
\end{tabular}
\caption{Twelve native public evaluations. Every refusal had tag
\code{ResidualBracketOutsideInterval}; every certification contained the
known exact root. The enclosure order and refinement budget were unchanged.}
\label{tab:matrix}
\end{table}

The matrix establishes noninterchangeability. F01 alone does not recover the
small affine displacement in B. F02 alone does not cure the below-one point
reduction in A. For C, resolving only one loss leaves the other active.
The combined result therefore supplies stronger evidence than testing each
patch only on its own favorable example.

The series constructions in the matrix were not regenerated for every variant;
the changes target certificate arithmetic, and the same stored inputs were
used. Reloading definitions in one kernel is useful for a controlled
comparison but does not substitute for clean-kernel installation and modular
loading checks before integration. Separate preliminary calls nevertheless
also reproduced the individual baseline and candidate mechanisms.

\subsection{Negative and fallback controls}
Additional executed native controls compared the baseline and combined
candidate. The nonlinear argument \code{1+x^2} on \code{{1,2}} produced an
identical enclosure. The positive argument \code{-x} on \code{{-2,-1}} also
produced an identical result. A logarithm of \code{1+x} over
\code{{-2,0}} retained an \code{IntervalDomain} refusal, as did a point
logarithm at zero. The value at one remained exactly \code{{0,0}}.

These controls address concrete risks of the edits: bypassing a domain check,
breaking the nonlinear fallback, or mishandling negative factors of a positive
argument. They are not an exhaustive equivalence check of all old and new
interval widths. A separate suspected lookup-default issue was dismissed by
a native control showing the expected shared lazy default; it is not included
as a third finding.

\subsection{Conditioning versus representation}
It would be inaccurate to say that $q\mapsto\log q$ is relatively
well-conditioned near one: its relative condition number is
$1/|\log q|$. The audit concerns an \emph{exact} rational $q$ whose small
displacement is known. With that displacement as input, the maps
\[
 \delta\mapsto\log(1+\delta),\qquad
 \delta\mapsto\log(1-\delta)
\]
have relative condition numbers tending to one as $\delta\to0$.
The corresponding inverse-root derivatives in the witnesses are also close
to one. The excessive uncertainty arises from changing that exact
representation into separately rounded or separately enclosed large terms.

A robust certificate engine should exploit exact structure when it is
available, while remaining conservative when it is not. The proposed edits
follow that principle without weakening any certified inclusion.

\section{Quantifying the avoidable work}
\subsection{A rational work model, not a benchmark}
The archive includes exact models of the rounded primitive and a separate
order-cost table. These calculations do not measure wall time, peak memory,
actual adaptive iterations, or total package cost. They show how two known
sources of uncertainty scale with $\delta=2^{-m}$.

For F01, define a simple comparison threshold by requiring
\begin{equation}\label{eq:proxy}
 T_n(2)\le\frac{2^{-m}}8.
\end{equation}
The smallest such $n\ge2$ is an exact rational work proxy for reducing the
near-$\log2$ tail to a fraction of the small target scale. It is not a
necessary and sufficient certificate-success condition. The actual
expression also contains the second tail, outward rounding, derivative
bounds, and interval geometry.

For F02, avoiding the particular rounded point interval
$[1,1+2^{1-B}]$ at $1+2^{-m}$ requires at least
\[
 B\ge m+1.
\]
With the driver's rule \eqref{eq:bits}, this gives the lower bound
\begin{equation}\label{eq:affineorder}
 n\ge \max\left(2,\left\lceil\frac{m+1}{4}\right\rceil-10\right).
\end{equation}
This lower bound only removes the identified grid loss; it does not prove
that all remaining certificate checks will succeed.

\begin{table}[htbp]
\centering
\begin{tabular}{@{}rrr@{}}
\toprule
$m$ in $\delta=2^{-m}$ & Least $n$ satisfying \eqref{eq:proxy} & Grid lower bound \eqref{eq:affineorder} \\
\midrule
100 & 31 & 16 \\
200 & 62 & 41 \\
500 & 156 & 116 \\
1000 & 314 & 241 \\
10000 & 3152 & 2491 \\
\bottomrule
\end{tabular}
\caption{Exact arithmetic work models, not kernel timings or measured retry
counts. The last row is a theoretical scale example, not an executed
10,000-bit public stress test.}
\end{table}

Since the driver caps $n$ at 2000, its largest $B$ is 8040. At $m=10000$,
early rounding of $1+2^{-m}$ still loses the small displacement even at that
cap. This identifies a precision obstacle in the original path, not a newly
executed cap-level failure. The narrow candidate preserves the exact
endpoints instead of attempting to hide the obstacle by raising a global cap.

For F01, \eqref{eq:proxy} implies an order growing asymptotically like
$m/(2\log_2 3)$ up to logarithmic corrections. With reciprocal reduction,
the series variable itself is of order $2^{-m}$, so a fixed number of terms
can already make the tail negligible relative to a fixed significant-bit
rounding budget.

\subsection{Representative exact-model widths}
At $m=200$, $n=2$, and $B=48$, the independent model gives
\begin{align*}
 \frac{\width(\I_{\rm original}(\log(1-\delta)))}{\delta}
   &\approx5.95\times10^{57},\\
 \frac{\width(\I_{\rm reciprocal}(\log(1-\delta)))}{\delta}
   &=2^{-47},\\
 \frac{\width(\I_{\rm early\ round}(\log(1+\delta)))}{\delta}
   &\approx1.14\times10^{46},\\
 \frac{\width(\I_{\rm affine\ fused}(\log(1+\delta)))}{\delta}
   &=2^{-48}.
\end{align*}
The CSV includes further scales from $m=20$ to $m=1000$. At smaller $m$, the
affine input may already be adequately representable, and F02 need not
improve the result. This is another reason not to advertise a universal
improvement factor.

The exact rational arithmetic still has input-size and intermediate-size
costs. In particular, a denominator $2^{m}$ has a binary representation of
length growing with $m$, and the finite rational powers in the unit-log
series involve larger integers. The claim is elimination of an artificial
need to increase \emph{series depth} with the small scale, not constant-time
arithmetic on arbitrarily large inputs.

\section{Wolfram 15 and Mathics3 implications}
\subsection{What was verified on Wolfram}
The main witnesses and edits were exercised on a real Wolfram 15.0.0 kernel,
loading the complete downloaded standalone. The experiment therefore went
beyond isolated helper transcription: public inverse objects were constructed
and passed through the actual certificate driver. Both candidate text anchors
were checked against the actual standalone, not merely against synthetic
fixtures.

The observations still have specific limits. They are selected remote
responses, not a complete MUnit log. The provided eight-test WLT file was not
executed as one \code{TestReport}. The candidate emitter was locally tested
on fixtures; equivalent substitutions were applied to the actual standalone
with Wolfram string operations. Neither successful mechanism implies the
canonical module was rebuilt and that its generated distribution was
revalidated.

\subsection{Why the source finding matters to Mathics}
The compatibility documentation targets Mathics3 10.0.1 with scanner 10.0.1,
SymPy 1.14.0, and Python 3.11. It explicitly distinguishes selected portable
acceptance from complete compatibility. The Mathics adapters isolate missing
or differing evaluator behavior in package-owned contexts rather than
patching global System functions.\cite{mathicscompat,mathicssite}

The logarithm certificate routines under review are shared exact-rational
code, rather than the separate Mathics numerical logarithm-splitting guard.
Their mathematical loss mechanisms do not depend on a machine floating-point
sign test. The candidate uses rational comparison, reciprocal construction,
interval negation, and the existing affine recognizer. These are appropriate
operations for a portable repair, but that source-level observation is not an
execution result.

There was no working Mathics kernel in this audit. Consequently the report
makes no claim that the candidate loads, returns early correctly through the
Mathics compatibility layer, meets its iteration budget, or produces exactly
the observed Wolfram matrix in Mathics. The supplied portable cases provide a
concrete acceptance target for that missing run.

\subsection{A focused cross-runtime acceptance obligation}
The proposed integration should execute the same A/B/C requests on both
runtimes and both supported layouts. Record whether each output is an
association or failure before reading numerical fields. On a successful
certificate, assert at least \code{Certified === True}, exact root
containment, and a center error bound below $\delta/8$. This prevents a
missing field or unevaluated expression from silently turning a numerical
comparison into a vacuous success.

The first acceptance run should retain the same order and zero-refinement
budget as the native evidence. Only afterward should adaptive defaults and
larger scales be measured. This separates local arithmetic behavior from the
already-existing driver's scheduling. The no-op exact-one, nonpositive-domain,
positive-negative-factor, and nonlinear fallback controls should travel with
the same patch.

The repository warns that a command-line expression parsed before package
loading can bind unfamiliar public symbols in the wrong context. The
supplied reproduction recipe therefore loads the package before the probe
as separate input expressions, and the probe uses fully qualified public
names. Its Mathics iteration-limit adjustment is explicit and matches the
repository's guidance; it is not an unannounced package-global change.
\cite{mathicscompat}

\section{E01: center-preserving logarithmic residuals}
\subsection{A narrower next step than a general expression optimizer}
F01 and F02 improve the logarithms supplied to the residual calculation.
There is a further, optional opportunity: sometimes the exact relationship
between the two logarithms in a residual can be retained instead of treating
them independently.

For strictly positive real $p$ and $q$,
\begin{equation}\label{eq:logratio}
 \log p-\log q=\log\left(\frac pq\right).
\end{equation}
This is an elementary real identity, not a new mathematical result. Its
application here should be a narrow, explicit certificate operation with
positivity obligations, not a global principal-log simplification rule.

For $p(x)=ax+b$ with rational coefficients and positive rational $q$, the
exact affine image of $p(x)/q$ can be passed to the logarithm routine. If an
exact rational center $c$ satisfies $ac+b=q$, the ratio at $c$ is exactly
one and its logarithm is exactly zero. Separate intervals for
$\log(ac+b)$ and $\log q$ usually retain nonzero uncertainty even when both
represent the same exact number.

\begin{proposition}[Centered affine log-ratio enclosure]
Let $a,b,q,\ell,r\in\Q$, $q>0$, $\ell\le r$, and
$ax+b>0$ throughout $[\ell,r]$. Define
\[
 J=\operatorname{sort}\left\{\frac{a\ell+b}{q},
                             \frac{ar+b}{q}\right\}.
\]
A containing logarithm interval for $J$ contains
$\log(ax+b)-\log q$ throughout the source interval. At a singleton $c$ with
$ac+b=q$, the existing exact-one rule produces the exact residual interval
$[0,0]$.
\end{proposition}
\begin{proof}
Division by positive $q$ preserves the positive affine image. Equation
\eqref{eq:logratio} holds on that domain, and monotonicity of the logarithm
then gives the enclosure. At the stated singleton the image is exactly
$\{1\}$, whose logarithm is exactly zero.
\end{proof}

\subsection{Executed primitive, unintegrated proposal}
With the combined candidate loaded, the following primitive was executed
natively and returned \code{{0,0}}:
\begin{lstlisting}
certLogExpression[(1+x)/(1+2^-200), x,
  {2^-200, 2^-200}, ctx]
\end{lstlisting}
The call used the actual private helper under its package context. It does
not establish that the public certificate dispatcher rewrites the original
residual to this ratio. That integration was neither implemented nor claimed.

The archive's \src{code/LogRatioPrototype.wl} provides explicitly named
rational and affine-ratio wrappers. The wrapper file itself was not executed
as a separate Wolfram or Mathics program. Its corresponding centered primitive
and the mathematical identity have the evidence stated above. It is a
prototype for a narrowly guarded operation, not an installed repair.

\subsection{Required guardrails}
A future dispatcher integration should derive the ratio from retained
positive-real data, not merely replace two syntactically matching
\code{Log} nodes. Native evaluation may rewrite a negative logarithm or
rearrange a rational target before the certificate routine sees it. A
syntactic match that accidentally broadens to complex inputs would lose the
real identity's hypotheses.

A zero point residual does not waive source-domain, continuity, derivative,
or root-containment obligations. An exact root does not by itself prove
uniqueness over the supplied interval or identify a global inverse branch.
The existing certificate checks must remain in charge of those claims.
The extension should preserve the original equation and identify any
transformed residual in its evidence, so an error bound is never silently
reinterpreted in another variable.

A sensible first implementation boundary is positive rational targets with
rational affine logarithm arguments only. That boundary is sufficient to
exercise exact centering, has an explicit proof, and avoids introducing a
new general logarithmic normalization subsystem. More general positive
expressions would require their own transformation and range contracts.

\section{Tests, artifacts, and reproducibility}
\subsection{Independent test population}
The independent Python run passed 24 unittest methods: five basic arithmetic
methods, thirteen logarithm/model methods, and six candidate-emitter fixture
methods. Some methods loop over many exact cases; those subcases are not
inflated into a package acceptance count. The actual output is retained in
\src{evidence/python-tests.txt}.

The exact model covers directed rounding for signed rationals, binary
exponents, interval arithmetic, baseline and repaired point logarithms,
affine transformations with positive and negative slopes, domain refusals,
reciprocity, and the explicit unit-tail formula. Near-one families cover
several small scales and series orders. A separate mean-value-radius model
checks a consequence of valid residual intervals; it is not presented as a
transcription of the entire certificate driver.

High-precision mpmath checks are independent numerical containment controls,
not proofs by themselves. One rational grid uses 180 decimal digits. The
near-one controls use \code{log1p} with 450 digits, including
$\delta=2^{-1000}$. Construction of the rational bounds never uses the
mpmath values or machine floating-point signs. The proofs in this article
supply the general inclusion arguments; numerical controls help detect
implementation mistakes in the independent model.

\subsection{Candidate emitter and application discipline}
The emitter accepts the canonical certificate module or a standalone file,
checks that each selected anchor occurs exactly once, and writes a new
candidate file. It refuses in-place replacement, an existing destination,
missing or repeated anchors, and repeat application. Its local tests use
synthetic anchor fixtures. This is useful fail-closed behavior, but not proof
of correct integration into an arbitrary later source revision.

For example, from the archive root:
\begin{lstlisting}
python code/emit_candidate.py \
  /path/to/repo/AsymptoticAnalysis.wl \
  /tmp/AsymptoticAnalysis.audit.wl --mode both
\end{lstlisting}
The other modes, \code{reciprocal} and \code{affine}, reproduce the separate
columns of the experiment. To integrate, apply the selected edits in a
staging checkout to \src{src/Kernel/InverseCertificates.wl}, regenerate the
standalone by the repository's maintained process, and execute the focused
acceptance there. This review did not perform that modular build.

\subsection{Native and portable reproduction}
Use separate input expressions after selecting a baseline or candidate path:
\begin{lstlisting}
Get["/tmp/AsymptoticAnalysis.audit.wl"];
Get["/path/to/audit/code/PortableLogCases.wl"];
\end{lstlisting}
The portable file contains three public witnesses and primitive observations.
Its expected baseline and candidate outcomes are the Wolfram-observed matrix,
not asserted Mathics results. The eight desired-contract MUnit cases can be
run on a native kernel after loading the combined candidate:
\begin{lstlisting}
TestReport["/path/to/audit/code/ReviewLogPrecision.wlt"]
\end{lstlisting}
That complete WLT invocation was not executed in this audit. The archive
retains this distinction even where its constituent cases correspond to
successful individual native probes.

From \src{code/}, the independent checks run with:
\begin{lstlisting}
python -m pip install -r requirements.txt
python -m unittest -v test_reference
\end{lstlisting}
The recorded Python environment was 3.13.5 with mpmath 1.3.0. This environment
is independent of the repository's documented Python 3.11 Mathics target;
passing the Python model does not validate that target interpreter.

\subsection{Evidence inventory}
\begin{longtable}{@{}p{2.2in}p{4.05in}@{}}
\toprule
Artifact & Meaning \\
\midrule
\endfirsthead
\toprule Artifact & Meaning \\ \midrule
\endhead
\path{evidence/scope.json} & Commit, runtime, completed work, and explicit
execution limits. \\
\path{evidence/findings.json} & Machine-readable descriptions of F01 and F02;
no inflated bug or test totals. \\
\path{evidence/native-matrix.json} & Selected transcription of the twelve
controlled public evaluations. Displayed decimal mantissas do not imply more
accuracy than the requested display precision. \\
\path{evidence/native-observations.json} & Individual primitive/public results,
negative controls, and the numerical-display warning. Not a raw MUnit log. \\
\path{evidence/python-tests.txt} & Actual output of the 24-method independent
model and fixture run. \\
\path{evidence/model-samples.csv} & Rational-model widths, rendered numerically
for comparison; not measured kernel timings. \\
\path{evidence/order-cost-model.json} & Exact tail-threshold and grid-order
models, including an explicitly theoretical large-scale row. \\
\path{evidence/novelty-crosswalk.md} & Nearest existing obligations and the
limits of the completed nonduplication screen. \\
\bottomrule
\end{longtable}

The source excerpts and anchors derive from the repository under its MIT
No Attribution license. The archive preserves that license and a provenance
notice. It does not redistribute the entire upstream package or its
historical review corpus.\cite{license}

\section{Recommended integration sequence and conclusion}
\subsection{Immediate scoped work}
First, integrate F01 and F02 as separately reviewable changes to the canonical
certificate module. Each has a small local mathematical obligation and an
individual witness. The combined gate must include witness C; otherwise a
translated version of the same logarithmic equation can remain obstructed
after a seemingly successful partial repair.

Second, run the supplied low-order cases in clean kernels on the intended
Wolfram and Mathics versions, in modular and standalone layouts. Preserve
exact source pinning and distinguish a script load failure, a structured
refusal, an unevaluated result, and a certified association. The present
native observations establish feasibility, not integration completion.

Third, measure only the changed paths at matched certificate goals. Include
admitted affine logarithms, nonaffine fallbacks, arguments away from one,
positive products with negative factors, and invalid domains. Report rational
operand sizes and achieved root widths alongside runtime. This measurement
should decide whether an extra failed affine probe is material; it should
not weaken the exact-endpoint invariant merely to reduce a dispatch count.

Finally, treat E01 as a separate feature experiment. The centered-ratio
primitive can preserve an exact point zero, but an integrated residual
transformation requires positivity evidence, equation provenance, and all
existing root checks. It should not delay the two smaller demonstrated
repairs.

\subsection{Final assessment}
The newly identified weakness is not that the package uses floating-point
arithmetic in its exact certificate proof. It does not in the primitives
examined here. The weakness is that sound enclosure operations can discard
valuable exact dependence and displacement information before the answer is
formed. Two different interfaces inside the logarithm path do so, and the
native matrix distinguishes them.

The repairs are correspondingly targeted: turn a below-one logarithm into
one well-scaled positive reciprocal logarithm, and move affine endpoint
rounding across the logarithm only when the exact affine range and positive
real domain justify it. The resulting public successes retain the safety
checks that caused the original refusals. This is a concrete improvement in
certificate reach and resource use, with explicit proofs and executable
regressions, rather than a proposal to accept weaker evidence.

The review's strongest conclusion is local and reproducible: on the pinned
standalone in Wolfram 15.0.0, three exact-root requests that refused at the
chosen fixed order certify after the combined changes, and every returned
interval contains the known exact root. Complete Mathics acceptance and
canonical-source integration remain to be demonstrated by the focused runs
provided with this article.

\clearpage
\begin{thebibliography}{99}
\small
\bibitem{source} AsymptoticAnalysis contributors.
\emph{Asymptotic repository}, revision
\path{efa1aeec4845a9c35e140963a0333d0c9ec33b05}, 10 September 2026.
\url{https://github.com/VladimirReshetnikov/Asymptotic/tree/efa1aeec4845a9c35e140963a0333d0c9ec33b05}.

\bibitem{certsource} AsymptoticAnalysis contributors.
\emph{InverseCertificates.wl}, pinned canonical source. In particular,
lines 8--18 (rounding), 54--72 (point logarithm), 80--130 (logarithm and affine
range), 257--313 (certificate attempt), and 365--490 (driver).
\url{https://github.com/VladimirReshetnikov/Asymptotic/blob/efa1aeec4845a9c35e140963a0333d0c9ec33b05/src/Kernel/InverseCertificates.wl}.

\bibitem{kernel} AsymptoticAnalysis contributors.
\emph{AsymptoticAnalysis.wl}, canonical kernel entry point and forward jet
arithmetic, same revision.
\url{https://github.com/VladimirReshetnikov/Asymptotic/blob/efa1aeec4845a9c35e140963a0333d0c9ec33b05/src/Kernel/AsymptoticAnalysis.wl}.

\bibitem{specialparam} AsymptoticAnalysis contributors.
\emph{ParameterizedSpecialFunctions.wl}, same revision.
\url{https://github.com/VladimirReshetnikov/Asymptotic/blob/efa1aeec4845a9c35e140963a0333d0c9ec33b05/src/Kernel/ParameterizedSpecialFunctions.wl}.

\bibitem{specialident} AsymptoticAnalysis contributors.
\emph{SpecialFunctionIdentities.wl}, same revision.
\url{https://github.com/VladimirReshetnikov/Asymptotic/blob/efa1aeec4845a9c35e140963a0333d0c9ec33b05/src/Kernel/SpecialFunctionIdentities.wl}.

\bibitem{reviews} AsymptoticAnalysis contributors.
\emph{Code review reports: index, six wave indexes, and retirement register},
same revision. Forty-five retained packages.
\url{https://github.com/VladimirReshetnikov/Asymptotic/blob/efa1aeec4845a9c35e140963a0333d0c9ec33b05/external-reports/code-review/README.md}.

\bibitem{status} AsymptoticAnalysis contributors.
\emph{Code review implementation status}, same revision. C19, C20, and P08
are the nearest maintained obligations discussed in the novelty crosswalk.
\url{https://github.com/VladimirReshetnikov/Asymptotic/blob/efa1aeec4845a9c35e140963a0333d0c9ec33b05/docs/development/CODE_REVIEW_STATUS.md}.

\bibitem{wavefive} AsymptoticAnalysis contributors.
\emph{Code review reports: wave 5}, same revision.
\url{https://github.com/VladimirReshetnikov/Asymptotic/blob/efa1aeec4845a9c35e140963a0333d0c9ec33b05/external-reports/code-review/wave-5/README.md}.

\bibitem{wavesix} AsymptoticAnalysis contributors.
\emph{Code review reports: wave 6}, same revision. Reports 46--55, including
its explicit overlap and evidence-boundary records.
\url{https://github.com/VladimirReshetnikov/Asymptotic/blob/efa1aeec4845a9c35e140963a0333d0c9ec33b05/external-reports/code-review/wave-6/README.md}.

\bibitem{mathicscompat} AsymptoticAnalysis contributors.
\emph{Running AsymptoticAnalysis in Mathics3}, pinned compatibility guide.
\url{https://github.com/VladimirReshetnikov/Asymptotic/blob/efa1aeec4845a9c35e140963a0333d0c9ec33b05/docs/Mathics/COMPATIBILITY.md}.

\bibitem{dlmf} NIST Digital Library of Mathematical Functions.
\emph{Section 4.6: Power Series}, especially equation 4.6.4 for the odd
logarithmic series. Consulted 10 September 2026.
\url{https://dlmf.nist.gov/4.6}.

\bibitem{logdoc} Wolfram Research.
\emph{Log: Wolfram Language documentation}. Background on logarithm
semantics and numerical evaluation; consulted 10 September 2026.
\url{https://reference.wolfram.com/language/ref/Log.html}.

\bibitem{mathicssite} Mathics3 project.
\emph{Mathics3 official website}. Consulted 10 September 2026.
The repository compatibility guide, not this general website, supplies the
specific tested dependency target used in this report.
\url{https://mathics.org/}.

\bibitem{license} AsymptoticAnalysis contributors.
\emph{MIT No Attribution License}, pinned repository license.
\url{https://github.com/VladimirReshetnikov/Asymptotic/blob/efa1aeec4845a9c35e140963a0333d0c9ec33b05/LICENSE}.
\end{thebibliography}
\end{document}
